% Standalone inclusion fragment; requires amsmath, amssymb, amsthm.
% Independent audit of the staged-ultrapower construction.
\section{Independence of progress: a coherent model of scalar computation}
\label{sec:progress-independence-audit}

The following argument takes place in classical metamathematics with choice.
Its conclusion is a countermodel to an intuitionistic calculus; using a
classical structure as an intuitionistic countermodel is sufficient.
No claim of a smaller complete axiomatization follows from this result.

\subsection{The semantic issue that must be addressed}

An arbitrary choice of an orbit for each \emph{written} update body is not
enough. Substituting terms for the passive parameters can change the written
body without changing its unary function. We therefore build an
\emph{extensional family} of scalar iterators. This is a stronger semantic
requirement than merely imposing each displayed computation equation, and
it prevents a gap in the native substitution lemma.

A body descriptor is a finite first-order term $h(u,\vec p)$ with one
distinguished state variable $u$ and a finite tuple of passive variables.
When an iterator for this descriptor is present, write its function symbol
as $I_h(n,a,\vec p)$. Its computation axioms are
\begin{align}
 I_h(0,a,\vec p)&=a,\label{eq:audit-c0}\\
 I_h(Sn,a,\vec p)&=h(I_h(n,a,\vec p),\vec p).
 \label{eq:audit-cs}
\end{align}
For \emph{every pair} of available descriptors, including two copies of the
same descriptor with different parameter values, require the first-order
schema
\begin{equation}
\begin{split}
 \forall\vec p\,\forall\vec q\bigl[
  (\forall u\ h(u,\vec p)=g(u,\vec q))\ \longrightarrow {}&\\[-0.6em]
  \forall n\,\forall a\ I_h(n,a,\vec p)=I_g(n,a,\vec q)
 \bigr].&
\end{split}
\tag{Ext}\label{eq:audit-ext}
\end{equation}
Thus an available iterator depends on the actual unary function and the
two numerical inputs, rather than on a particular presentation of that
function. Parameter repetition, omission, and reordering are included.

\subsection{The orbit supplied by one ultrapower}

Let $A$ be any nonempty structure, let $U$ be a nonprincipal ultrafilter on
$\mathbb N$, and put $B=A^{\mathbb N}/U$. Fix a unary term function of $B$,
\[
 f(v)=h^B(v,\vec p),\qquad \vec p=[\vec p_i],
\]
and an element $a=[a_i]$ of $B$. Write
$f_i(v)=h^A(v,\vec p_i)$. For each ordinary integer $j\in\mathbb Z$, set
\[
 x_j=\bigl[f_i^{\,\max(0,i+j)}(a_i)\bigr].
\]
For all sufficiently large $i$,
$\max(0,i+j+1)=\max(0,i+j)+1$. Hence
\[
                         f(x_j)=x_{j+1}
                         \quad(j\in\mathbb Z).
\]
Every parameter-definable unary term function in this ultrapower therefore
has a full two-sided orbit. The iterations in the coordinates are finite
external iterations; no arithmetic induction axiom in $A$ is being used.
The displayed sequence need not pass through $a$, and no such condition is
required on a successor component that does not contain zero.

\subsection{Extension of a coherent partial family}

Suppose $A$ has successor injectivity, zero outside the image of successor,
zero/successor coverage, and no finite successor cycles. Suppose it also has
a collection of iterator symbols satisfying
\eqref{eq:audit-c0}, \eqref{eq:audit-cs}, and all instances of
\eqref{eq:audit-ext} for that collection.

Take the ultrapower $B$ in the language containing \emph{all} those symbols.
All the stated first-order conditions transfer by \L o\'s's theorem.
In particular, the old iterator values are retained, and their coherence
holds also for parameter tuples newly present in the ultrapower.

The successor graph of $B$ consists of the standard ray
$0,S0,S^2 0,\ldots$ and components isomorphic to $\mathbb Z$.
Indeed, every nonzero point has a unique predecessor. A point outside the
standard ray cannot reach zero after finitely many predecessor steps;
injectivity and the absence of finite cycles ensure that all steps in its
two-sided component are distinct. Choose an origin in each such component.

Now add iterators for any set of descriptors in the \emph{old language}.
Partition their parameter instances by the actual unary functions
$f:B\to B$ that they define.

If a function $f$ already has an iterator from the old family, use that
iterator at every counter and every seed. This prescription is well-defined:
if $f=h^B(-,\vec p)=g^B(-,\vec q)$ for two old descriptors, Ext says that
their iterator values agree for all counters and seeds.

For a function not represented in the old family, make one simultaneous
choice for its entire equivalence class. At a standard counter $S^k0$ put
\[
                         I_f(S^k0,a)=f^k(a).
\]
On a nonstandard successor component with chosen origin $b$, choose a
two-sided orbit $(x_j)_{j\in\mathbb Z}$ supplied by the preceding lemma,
and set
\[
                         I_f(S^j b,a)=x_j.
\]
Here negative powers mean the unique iterated predecessor on that component.
Choices may be made separately for each $f$, seed $a$, and component.
The set-theoretic axiom of choice suffices for all these choices.

Define every newly added $I_h(n,a,\vec p)$ by the chosen iterator for the
actual function $h^B(-,\vec p)$. Its two computation equations hold by
construction. Ext holds for every pair of old or new symbols, because
equal unary functions were assigned the same entire iterator. The expansion
changes no old function or old element.

\subsection{A seed model in which identity iteration is not constant}

Start with the disjoint union
\[
                         A_0=\mathbb N\ \sqcup\ \mathbb Z.
\]
Zero belongs to the first summand; successor is ordinary successor on each
summand. Let $1=S0$. Interpret the iterator $J=I_u$ of the identity update by
\[
 J(n,a)=
 \begin{cases}
 a,&n\in\mathbb N,\\
 1,&n\in\mathbb Z.
 \end{cases}
\]
Then $J(0,a)=a$ and $J(Sn,a)=J(n,a)$. All the required successor conditions
hold. The initial extensionality condition for the sole descriptor $u$ is
automatic. For any chosen $c$ in the $\mathbb Z$ summand,
\[
                           J(c,0)=1\ne0.
\]
No predicate distinguishing the two summands is required in the language.

Iterate the extension construction through $\omega$ stages. At each stage,
take an ultrapower in the entire current language and add iterators for all
finite unary term descriptors of that language. Identify the preceding
structure with its diagonal elementary image before adding the symbols.
The languages and structures can then be regarded as increasing chains.
The reduct at every later stage is an elementary extension in each earlier
language. Their union $A_\omega$ retains all computation equations, all
cross-extensionality schemas, all successor conditions, and the distinguished
equation $J(c,0)=1\ne0$.

To justify the union step explicitly, each fixed old-language sentence is
handled by the elementary-chain theorem on the reducts to that language.
Every computation or Ext instance uses only finitely many symbols and was
imposed at a finite stage, so it holds in the union. Equivalently, Ext can
be checked directly: if two final unary functions agree, they agree on
every sufficiently late stage containing their parameters; applying Ext at
a stage also containing the proposed counter and seed gives equality of
the two iterator values there and hence in the union.

\subsection{Returning to native scalar syntax}

Every finite native term
\[
 t::=x\mid0\mid St\mid I(n,a;u.h)
\]
has an interpretation using finitely many staged symbols. Expand its inner
iterators first. The translated update body is then a finite term of some
stage language, so its descriptor has an iterator at the next stage.
Parameters are ordinary numerical variables; there is no function sort and
no internal universal evaluator.

More intrinsically, for any unary function definable by a finite native
body with parameters in $A_\omega$, choose a staged descriptor presenting
it and use its iterator. Ext makes this choice independent of the
descriptor. It follows by structural induction on native terms that
capture-avoiding substitution has the usual semantics:
\[
 \llbracket t[s/x]\rrbracket_\rho
 =\llbracket t\rrbracket_{\rho[x\mapsto\llbracket s\rrbracket_\rho]}.
\]
At an iterator, the induction hypothesis identifies its counter, seed, and
pointwise update function; Ext identifies the resulting iterators. This is
the exact point at which mere independent choices for written bodies would
have been insufficient.

\subsection{Soundness without progress and the refuted instance}

Interpret native equality as actual equality, and use the ordinary readings
\[
 Q_y B=B(0)\to\forall y\,B(y),\qquad
 D_{x,y}B=\exists x\bigl(B(x,0)\to\forall y\,B(x,y)\bigr).
\]
Hypothesis, equality transport, the two computations, closure elimination,
and witness elimination are sound. Case introduction is sound because every
element is zero or a successor: from $B(t,0)$ and all successor instances
$B(t,Sz)$ one obtains every instance. The usual freshness conditions ensure
that $t$ and the retained assumptions do not depend on the fresh $z$.
All these arguments are ordinary first-order soundness arguments in the
single structure $A_\omega$.

Let
\[
                          E(y)\;:\equiv\;I(y,0;u.u)=0.
\]
The zero computation proves $E(0)$. The successor computation and ordinary
derived equality reasoning prove
\[
                          E(z)\ \vdash\ E(Sz)
\]
without any induction, since
$I(Sz,0;u.u)=I(z,0;u.u)=0$ under the displayed hypothesis.
Nevertheless $E(c)$ is false in $A_\omega$. Consequently $Q_yE(y)$ is
false there and is not derivable by the kernel without progress.
Thus QI by progress is not admissible over the other seven rules.
With an unused existential witness, the same model also refutes the
witnessed progress conclusion $D_{x,y}E(y)$.

\paragraph{Audit status and limits.}
This is a written model construction, not a proof-assistant verification.
The essential requirements are simultaneous extensional choice for actual
unary functions; universal cross-extensionality for all parameter tuples;
ultrapowers in the full current language; and adding only prior-language
body descriptors at a given stage. Omitting any of those requirements would
leave a genuine coherence or circularity gap. The result proves independence
of progress. It does not prove a complexity lower bound for other
presentations, and it does not satisfy a requirement for a smaller complete
proof checker by itself.
